\section{Positive logarithmic energy on an interval}\label{sec:p0-analytic}

This section records the analytic passage for the positive P0 determinant,
including the dependence on the base.  The finite determinant statement is
for integers $q$ with $2\le q\le11$; a real parameter is used only in the
continuum interpolation.  All measures, intervals, and the contour below are
fixed as $q$ varies.  The two endpoint reference measures are the exact data
in Tables~\ref{tab:q2data} and~\ref{tab:q11data}.

Put $m\in4\mathbb N$, $m\ge8$, $g=m/2$,
\[
 I_m=\{m/4,\ldots,5m/4-1\},\qquad
 J_m=\{0,\ldots,2m-1\},\qquad E=[1/4,5/4],
\]
and fix
\[
 \Gamma=\{\xi+A\cos u+iB\sin u:0\le u<2\pi\},\quad
 (\xi,A,B)=\left(\frac{11281}{10000},\frac{12783}{10000},
                                      \frac{229}{200}\right).
\]
Write $\Delta(X)=\prod_{a<b}(x_b-x_a)$ and $K(x,y)=(x+y)^{-1}$.
The source function and the frequency weight at base $q>1$ are
\[
 F_{g,q}(z)=q^z\binom{z+g}{g},\qquad
 U_i(q)=q^i\binom{i+g}{g},\qquad
 w_{r,q}=\binom{r-1}{g}q^{-r}\quad(r>g).
\]
For integer $q$, let $G_{m,q}$ be the P0 determinant of
Section~\ref{sec:p0}, with $a=q$.  The diagonal contour identity is
\begin{equation}\label{eq:p0-contour-q}
 C^{(g)}_{ij}(q)=-\oint_{m\Gamma}
       \frac{F_{g,q}(z)}{(z-i)(z-j)}\frac{dz}{2\pi i}
             +(\log q)U_i(q)\,\mathbf1_{i=j}.
\end{equation}
For $i\ne j$ this is the divided-difference residue identity; for $i=j$ it
follows by differentiating $F_{g,q}$, since
\[
 F_{g,q}'(i)=U_i(q)\left(\log q+\sum_{s=1}^g\frac1{i+s}\right).
\]
Expanding the rational moment expression in the frequency basis and applying
Cauchy--Binet~\cite{Gantmacher} gives an absolutely convergent sum over
$R\subset\{g+1,g+2,\ldots\}$, $|R|=m$, with weight
$\prod_{r\in R}w_{r,q}$.  After choosing $k$ diagonal rows $S\subset I_m$
and putting $d=m-k$, its absolute integrand is
\begin{equation}\label{eq:p0-q-integrand}
 \begin{aligned}
 Q_{m,q}(S,R,Z)={}&\left(\prod_{i\in S}(\log q)U_i(q)\right)
   \left(\prod_{r\in R}w_{r,q}\right)
   \left(\prod_{z\in Z}|F_{g,q}(z)|\right)\\
 &\times\left|\det K_{I_m\setminus S,-Z}\det K_{I_m,R}
                \det K_{(-Z,R),J_m\setminus S}\right|.
 \end{aligned}
\end{equation}
where $Z=(z_1,\ldots,z_d)\in(m\Gamma)^d$.  The integration measure is
$\prod |dz|/(2\pi)$ and the symmetrization factor is $1/d!$.  The three
Cauchy factors have, after restoring the deleted fixed-grid factors, the
same three Vandermonde interactions as at $q=2$: contour-contour, pinned-
pinned, and frequency-frequency each have coefficient one; pinned-contour
has coefficient two; pinned-frequency and contour-frequency each have
 coefficient one.  In particular, no factor depending on $q$ enters the
 quadratic interaction energy.

\begin{proof}[Cauchy--Binet energy expansion~\cite{Gantmacher}]
The rational moment expression is obtained by expanding the two Cauchy kernels
geometrically and using the exact generating identity
\[
 \sum_{r\ge g+1}\binom{r-1}{g}(t/q)^r
       =\left(\frac{t}{q-t}\right)^{g+1},\qquad 0\le t<q.
\]
Inserting these moments and applying
\[
 \det(XY)=\sum_{|R|=m}\det X_{:,R}\det Y_{R,:}
\]
gives the sum over the full frequency sets $R$ stated above.  Splitting the
diagonal term in \eqref{eq:p0-contour-q} row by row chooses
$S\subset I_m$; the remaining $d=m-|S|$ rows are represented by contour
variables $Z$, with the factor $1/d!$ from their antisymmetrization.
For any two equal-size sets, the Cauchy determinant identity~\cite{Krattenthaler}
\[
 \det\left(\frac1{x_i+y_j}\right)_{i,j=1}^r
 =\frac{\Delta(x)\Delta(y)}{\prod_{i,j}(x_i+y_j)}
\]
converts each of the three kernel minors into a Vandermonde quotient.
Restoring the fixed-grid Vandermonde factors leaves contour-contour,
pinned-pinned, and frequency-frequency interactions with coefficient one,
pinned-contour with coefficient two, and the other two mixed interactions
with coefficient one.  Taking absolute values gives exactly
\eqref{eq:p0-q-integrand}; all signs and $q$-powers are one-body terms.
For fixed $m$, the Cauchy minors have at most polynomial growth in the
frequencies and $w_{r,q}=O_q(r^gq^{-r})$.  The absolute frequency sum
therefore converges, and increasing finite truncations converge to its
full absolute envelope.
\end{proof}

For the normalization, the signed Cauchy--Binet expansion~\cite{Gantmacher}
can be written
explicitly. Put $n=m/4$, $Y=S\uplus Z$, and
\[
 \begin{aligned}
 D_m(a)&=(a-n)!(n+m-1-a)!a!(2m-1-a)!,\\
 P_m(z)&=\prod_{a\in I_m\uplus J_m}(z-a),\qquad
 Q_m^*(r)=\prod_{a\in I_m\uplus J_m}(r+a).
 \end{aligned}
\]
Then
\begin{align*}
G_{m,q}={}&\Delta(I_m)^2\Delta(J_m)
 \sum_{S\subset I_m}\frac{(-1)^{|S|(n+1)}(\log q)^{|S|}}{d!}
 \sum_{|R|=m}
 \left(\prod_{a\in S}\frac{U_a(q)}{D_m(a)}\right)
 \left(\prod_{r\in R}w_{r,q}\right)\\
&\quad\times\oint_{(m\Gamma)^d}
 \frac{\Delta(Y)^2\Delta(R)^2\prod_{y\in Y,r\in R}(y+r)}
      {\prod_{z\in Z}P_m(z)\prod_{r\in R}Q_m^*(r)}
      \prod_{z\in Z}\left(F_{g,q}(z)\frac{dz}{2\pi i}\right).
\end{align*}
This follows by applying Cauchy--Binet~\cite{Gantmacher} to the two Cauchy
moment factors,
then expanding the diagonal summands in the contour formula row by row.
For every equal-size pair, use the Cauchy determinant identity~\cite{Krattenthaler}
\[
 \det\left((x_i+y_j)^{-1}\right)_{i,j=1}^r
 =\frac{\Delta(x)\Delta(y)}{\prod_{i,j}(x_i+y_j)}.
\]
Restoring the deleted fixed-grid factors uses
\[
 \Delta(I_m\setminus S)\Delta(J_m\setminus S)
 =\frac{\Delta(I_m)\Delta(J_m)\Delta(S)^2}
        {\prod_{a\in S}D_m(a)}.
\]
Each pinned index is therefore counted twice in a contour denominator and
once in a frequency denominator. The three Cauchy factors give
contour-contour, pinned-pinned, and frequency-frequency interactions with
coefficient one, pinned-contour with coefficient two, and each remaining
mixed interaction with coefficient one. Taking absolute values and scaling
all positions by $m$ gives \eqref{eq:p0-q-integrand} and
\eqref{eq:p0-q-energy}. The factors $q^a$, $q^{-r}$, and $q^z$ contribute
only the one-body fields $\lambda\Re z$ and $-\lambda t$ in
\eqref{eq:p0-q-fields}; the quadratic interaction coefficients are
independent of $q$.

The fixed-grid normalization is also explicit.  Consecutive integer lists
satisfy $\Delta(I_m)=\prod_{j=1}^{m-1}j!$ and
$\Delta(J_m)=\prod_{j=1}^{2m-1}j!$.  Stirling's formula~\cite{WhittakerWatson}
summed over these
products gives
\[
 \log(\Delta(I_m)^2\Delta(J_m))
 =3m^2\log m+(2\log2-9/2)m^2+O(m\log m).
\]
Equivalently, the fixed Vandermonde contributes $C_0=9/2-2\log2$ to
the negative logarithmic energy.  The scaled Cauchy quotients cancel all
$m^2\log m$ terms; their remaining homogeneous factor is
$m^{-2m-2d}$ before the contour differentials and $m^{-2m-d}$ after them.
This accounts for the full quadratic normalization in
\eqref{eq:p0-q-energy}.

\subsection{The fixed variational problem and its affine parameter}

Let $\eta=\kappa+\nu$, where
\[
 0\le\kappa\le1_E\,dx,\quad \nu\ge0,\quad
 \operatorname{supp}\nu\subset\Gamma,\quad |\kappa|+|\nu|=1,
\qquad 0\le\mu\le1_{[1/2,20]}dt,\quad |\mu|=1.
\]
Write $\mu^-=(-\operatorname{id})_\#\mu$ for reflection onto the negative
real axis.
The set $\mathcal A$ of such pairs $(\eta,\mu)$ is independent of $q$.
Define
\[
 A_0(z)=\int_E\log|z-x|\,dx+\int_0^2\log|z-x|\,dx,
 \qquad
 B_0(t)=\int_E\log(t+x)\,dx+\int_0^2\log(t+x)\,dx,
\]
\begin{align*}
 \phi_2(z)&=(\log2)\Re z+\int_0^{1/2}\log|z+s|\,ds
                              +\frac12+\frac12\log2,\\
 \psi_2(t)&=t\log t-(t-1/2)\log(t-1/2)
                         +\frac12\log2-t\log2.
\end{align*}
For measures on the corresponding reflected supports write
\[
 I(\sigma,\tau)=-\iint\log|x-y|\,d\sigma(x)d\tau(y).
\]
The energy is
\begin{align}
 \mathcal E_{q}(\eta,\mu)={}&C_0+
     \int(A_0-\phi_q)\,d\eta+\int(B_0-\psi_q)\,d\mu \\
 &+I(\eta)+I(\mu)+I(\eta,\mu^-),
 \label{eq:p0-q-energy}
\end{align}
Here $I$ is positive on zero-mass variations, while $Q_{\log}$ below uses
the opposite, positive-logarithm convention.  Also
$C_0=9/2-2\log2$ and, on writing $q=2e^\lambda$,
\begin{equation}\label{eq:p0-q-fields}
 \phi_q(z)=\phi_2(z)+\lambda\Re z,
 \qquad \psi_q(t)=\psi_2(t)-\lambda t.
\end{equation}
The logarithmic integrals are understood in the following regularized sense:
truncate at the logarithmic singularity, integrate the resulting bounded
functions, and take the limit after replacing the coincident-cell term by
the continuous primitive \(u(\log|u|-1)\), with value \(0\) at \(u=0\).
The measures used below have bounded densities and compact support, so this
limit is finite.  Substitution into
\eqref{eq:p0-q-energy} shows, for every fixed admissible pair,
\begin{equation}\label{eq:p0-q-affine-energy}
 \mathcal E_\lambda(\eta,\mu)=\mathcal E_0(\eta,\mu)
       +\lambda\left(-\int\Re z\,d\eta(z)+\int t\,d\mu(t)\right).
\end{equation}
Consequently
\begin{equation}\label{eq:p0-q-alpha}
 \alpha(\lambda)=\inf_{(\eta,\mu)\in\mathcal A}
                    \mathcal E_\lambda(\eta,\mu)
\end{equation}
is concave.  Indeed, for $0\le\theta\le1$, each fixed-pair energy at
$\theta\lambda_0+(1-\theta)\lambda_1$ is the corresponding convex
combination of its endpoint energies, and taking an infimum gives
\[
 \alpha(\theta\lambda_0+(1-\theta)\lambda_1)
       \ge\theta\alpha(\lambda_0)+(1-\theta)\alpha(\lambda_1).
\]
Apply the shared quadratic-duality identity of
Section~\ref{sec:certificate-algorithm} with quadratic part
$Q_{\rm P0}(\eta,\mu)=I(\eta)+I(\mu)+I(\eta,\mu^-)$.  Fix an admissible
reference $(\eta_0,\mu_0)$ and put
\begin{align*}
 Q_{\log}={}&\iint\log|z-w|\,d\eta_0(z)d\eta_0(w)
 +\iint\log|t-s|\,d\mu_0(t)d\mu_0(s)\\
 &+\iint\log|z+t|\,d\eta_0(z)d\mu_0(t).
\end{align*}
\[
Q_{\rm P0}(\eta_0,\mu_0)=-Q_{\log},\qquad
W_\eta=A_0-\phi_q-2L_{\eta_0}-L_{\mu_0^-},\qquad
 W_\mu=B_0-\psi_q-2L_{\mu_0}-L_{\eta_0}(-t).
\]
If $W_\eta\ge h_\eta$ on $\Gamma$, set
\[
 P=\int_E(h_\eta-W_\eta)_+\,dx,\qquad
 T=\int_{1/2}^{20}(h_\mu-W_\mu)_+\,dt.
\]
Then
\begin{equation}\label{eq:p0-q-dual}
 \alpha(\lambda)\ge C_0+Q_{\log}+h_\eta+h_\mu-P-T.
\end{equation}
For $\sigma=\eta-\eta_0$ and $\tau=\mu^- -\mu_0^-$, the shared identity
has remainder
\[
 I(\sigma)+I(\tau)+I(\sigma,\tau)
 =\tfrac12 I(\sigma)+\tfrac12 I(\tau)
                         +\tfrac12 I(\sigma+\tau)\ge0,
\]
where $I(\rho)=-\iint\log|x-y|\,d\rho(x)d\rho(y)$ is nonnegative for
zero-mass signed measures.  The density constraints then give
$\int W_\eta d\eta\ge h_\eta-P$ and
$\int W_\mu d\mu\ge h_\mu-T$, proving \eqref{eq:p0-q-dual}.

\subsection{Finite configurations and cell estimates}

For a selected pinned set $S\subset I_m$ and a frequency set $R$, define
\begin{align*}
 \kappa_S&=\sum_{i\in S}\mathbf1_{[i/m,(i+1)/m)}dx,
 &\mu_R&=\sum_{r\in R}\mathbf1_{[(r-1)/m,r/m)}dt,\\
 \nu_\delta&=\frac1m\sum_{a=1}^{d}\sigma_{z_a/m,\delta},
 &\delta&=m^{-2}.
\end{align*}
Here $\sigma_{z,\delta}$ is probability arclength on the radius-$\delta$
circle centered at $z$.
Put $\zeta_a=z_a/m\in\Gamma$.  Thus all logarithms in the continuum
energy are in scaled coordinates; the powers of $m$ arising from the change
from $z_a$ to $\zeta_a$ have already been included in the fixed
Vandermonde normalization in the cell estimate.
The first two measures have density at most one,
masses $|S|/m$ and one, and the first frequency cell begins at $1/2$.
Circle averaging gives
\begin{equation}\label{eq:p0-circle}
 \sum_{a\ne b}\log|\zeta_a-\zeta_b|
 \le m^2\iint\log|z-w|\,d\nu_\delta(z)d\nu_\delta(w)
                                          +2d\log m.
\end{equation}
The fixed-grid and reflected-frequency logarithmic potentials are harmonic
on the smoothing disks.  The contour field is subharmonic, including across
the zero segment of $F_{g,q}$.  Thus smoothing has no additional field
error; the explicit $2d\log m$ is the full smoothing loss.

The exact real-cell estimates used in the envelope are listed here to make
the finite-size accounting reproducible.  For $n$ distinct lattice sites,
the cell self-energy inequality is
\[
 \sum_{a\ne b}\log|x_a-x_b|
 \le m^2\iint\log|x-y|\,d\rho(x)d\rho(y)
                                  +n(\log m+\log(2\pi)),
\]
where $\rho$ is the union of their length-$1/m$ cells.  To see the constant,
write
\[
 e_d=\log d-\int_0^1\!\int_0^1\log|d+u-v|\,du\,dv.
\]
Then $e_1=3/2-2\log2$ and for $d\ge2$,
\[
 e_d=\sum_{j\ge1}\frac1{j(2j+1)(2j+2)d^{2j}}
       \le\frac1{12d(d-1)},\qquad
 \sum_{d\ge1}e_d=\tfrac12\log(2\pi)-\tfrac34.
\]
Indeed, for $d\ge2$ expand
$\log(1+(u-v)/d)$ in its absolutely convergent power series and integrate
term by term; the even moments
$\int_0^1\!\int_0^1(u-v)^{2j}\,du\,dv=2/((2j+1)(2j+2))$
give the displayed series, while the $d=1$ integral is elementary.  Summing
the same identity from $d=1$ to $M$, exchanging the positive series, and
using Stirling's Wallis limit~\cite{WhittakerWatson}
$\log M!=(M+\tfrac12)\log M-M+\tfrac12\log(2\pi)+o(1)$
gives the displayed infinite sum.  Thus no numerical cell constant is being
imported.
There are at most two sites at each distance from a fixed site.  This proves
the displayed inequality by summing the self-cell and pair-cell errors.
Applied to the selected pinned and frequency sites and to the fixed
Vandermonde, it contributes $4m(\log m+\log(2\pi))$.

The remaining one-body and mixed-cell bounds are as follows.  Binomial
entropy gives
\[
 \log U_i(q)\le m\phi_q(i/m),\qquad
 \log w_{r,q}\le m\psi_q(r/m).
\]
The exact cell corrections from the affine parameter are
$-\lambda m/2$ for frequencies and $-\lambda|S|/2$ for pinned cells.
For the base-$2$ frequency field,
\[
 \psi_2'(t)=\log\frac{t}{t-1/2}-\log2,\qquad
 \int_{1/2}^{1}(\psi_2'(t))_+\,dt=\tfrac12\log2.
\]
Summing endpoint-to-cell differences over disjoint frequency cells costs
at most $(m/2)\log2$.  Monotone quadrature for the fixed-grid field gives,
for $t\ge1/2$,
\[
 0\le mB_0(t)-\sum_{i\in I_m}\log(t+i/m)
                    -\sum_{j\in J_m}\log(t+j/m)
 \le\log(35/3).
\]
Thus the frequency one-body cost is at most
$m(\log(35/3)+\tfrac12\log2)-\lambda m/2$.

For the pinned denominator, take $N$ consecutive sites starting at $a$ and
containing $i$, with $n_L=i-a$ and $n_R=a+N-1-i$.  Comparing the omitted
site product on the two sides of $i$ with its cell integral gives
\[
 \log(n_L!n_R!)-(N-1)\log m
 -m\int_{a/m}^{(a+N)/m}\log|i/m-y|\,dy
 \ge\log\frac{m}{n_R+1}.
\]
This follows from monotone logarithmic sums and
$n_R\log(1+1/n_R)\le1$; the zero case is direct.  The $I_m$ contribution
is nonnegative and the $J_m$ contribution is at least $-\log2$.
Direct integration of the explicit derivatives gives
$\int_E| (\phi_2-A_0)'|\le\log6+2\log2+\log7=\log168$.
Consequently pinned quadrature costs at most $m\log336$.  Its affine
correction is $-\lambda|S|/2$, since the linear-field integral on each
cell is centered one half-cell to the right of the selected lattice point.
For an exact bookkeeping check, $q/2-\log q$ is increasing and positive on
$[2,\infty)$, hence $(\log\log q)_+\le\lambda$.  Therefore the diagonal
and cell-shift contribution satisfies
\[
 |S|(\log\log q)_+-\frac\lambda2(|S|+m)
 \le-\frac\lambda2(m-|S|)\le0.
\]
The displayed pointwise estimate keeps the positive term and drops these
favorable shifts, so it remains valid without relying on this cancellation.

The ellipse clearances are
\[
 d_+=\operatorname{dist}(\Gamma,[0,2])=751/5000,\qquad
 d_-=\operatorname{dist}(\Gamma,(-\infty,-1/2])=1749/5000.
\]
They follow by minimizing the squared distance as a quadratic in
$\cos u$ on the ellipse parametrization.  The pinned-contour displacement
costs $m/d_+$, the fixed contour denominators cost $3m/(2d_+)$, and the
contour-frequency displacement costs $m/(2d_-)$.  The pinned-frequency
cell displacement is centered; Taylor's inequality, with
$i/m+r/m\ge3/4$ and cell width $1/m\le1/4$, bounds its total loss by $1/3$.
Together these contributions are at most
\[
 m\left(\frac5{2d_+}+\frac1{2d_-}\right)+\frac13.
\]
Set $T_\delta=\{z:\operatorname{dist}(z,\Gamma)\le\delta\}$, with the tube
radius $\delta$ used in the transfer argument below.
For the finite contour product the required inequality is one-sided.  Fix
$z\in T_\delta$, put $\zeta=mz$, $r=-\Re\zeta$, $b=\Im\zeta$, and
$f(u)=\log|\zeta+u|$.  Then
\[
 f'(u)=\frac{u-r}{(u-r)^2+b^2},
\]
so $f$ decreases to $r$ and increases after $r$.  Put $H=\max(f,0)$.
The negative part of $f$ is supported near $r$; after translating $r$ and
using $\int_{-1}^{1}(\log|u|)_-\,du=2$, its integral over every unit cell
meeting $r$ is at most $2$.  On each of the two monotone pieces the
  right-endpoint quadrature error for $H$ is nonpositive on the decreasing
  piece and nonnegative on the increasing piece.  The positive contribution
  is bounded by the variation on the increasing piece.  The sampled endpoint
  values are bounded by
$\log(4(g+1))$ near the turning cell and by
$\log(6(g+1))$ at the right endpoint, so this variation is at most
$\log(6(g+1))$.  Thus the complete discrepancy satisfies
\[
 \sum_{k=0}^{g-1}f(k+1)-\int_0^g f(u)\,du
 \le 2+\log(6(g+1)).
\]
This is the bounded-variation estimate for the truncated function $H$; it
also covers the cases where the zero of $f$ lies outside $[0,g]$.  We retain
the looser uniform quantity
\[
 \beta_m=4\log(g+1)+9+10\log2
\]
below.  Restoring the factors $\log m$ gives the cell integral in the
contour-product estimate.
Thus, for every $z\in T_\delta$,
\begin{equation}\label{eq:p0-q-contour-product}
 \log|F_{g,q}(mz)|\le m\phi_q(z)+\beta_m.
\end{equation}
The scalar factor $q^{mz}$ is exactly the affine field term and adds no
discrepancy.  No lower bound for this discrepancy is used.
Since $g+1\le m$ for $m\ge8$, the contour
and circle errors together contribute $2m\log m+m\beta_m$.

Combining the preceding terms with the fixed Vandermonde normalization and
discarding favorable powers of $m$ gives, for every configuration,
\[
 \log Q_{m,q}(S,R,Z)
 \le -m^2\mathcal E_q(\kappa_S+\nu_\delta,\mu_R)
   +10m\log m+c_*m+\frac13
   +|S|\log\log q,
\]
where the real-cell constant is
\begin{align*}
c_*={}&4\log(2\pi)+9+\frac{21}{2}\log2+\log336+\log(35/3)
                         +\frac5{2d_+}+\frac1{2d_-}\\
 <{}&49.978.
\end{align*}
The displayed logarithm bounds can be checked with the rational series
\[
 \log x=2\sum_{j=0}^{N}\frac{u^{2j+1}}{2j+1}+R_N,\qquad
 0<R_N<\frac{2u^{2N+3}}{(2N+3)(1-u^2)},\quad
 u=\frac{x-1}{x+1}.
\]
For $x=710/113$ use $N=50$; for
$x=2,21/16,35/24,39/32$ use $N=16$.  The rational series gives
\[
 \log(710/113)<1.83787716,\quad \log2<0.69314719,
 \quad \log(21/16)<0.27193372,
\]
\[
 \log(35/24)<0.37729424,\qquad \log(39/32)<0.19782575.
\]
The classical rational bound $\pi<355/113$ therefore gives
$\log(2\pi)<1.83787716$; hence
$\log336<8(0.69314719)+0.27193372$ and
$\log(35/3)<3(0.69314719)+0.37729424$.
Substitution of the exact rational $d_+=751/5000$ and
$d_-=1749/5000$ gives $c_*<49.977264<49.978$ and
$c_*+\log39+1+A<55.920<56$.
This is the base-$2$ cell constant; the $q$-dependent cell shifts above are
favorable.  The only $q$-dependent finite factor not absorbed by the affine
energy is the diagonal multiplier $(\log q)^{|S|}$.  In the pointwise
envelope it appears as the exact term $|S|\log\log q$.  At $q=2$ this term
is negative, so the base-$2$ estimate omits that favorable factor; whenever an
upper bound uniform in $q$ is taken below, it is replaced by
$|S|(\log\log q)_+$.  Uniformly
for $2\le q\le11$,
\[
 \log q<3,\qquad 0<\log q<3<5,\qquad \log(\log q)<\log3<2,
\]
so $|S|(\log\log q)_+<2m$.  For all bases we therefore replace $56m$ by at
most $58m$; this
finite correction is $O(m)$ and does not change the limiting energy rate.

\subsection{Tube transfer, frequency tail, and summation}

Let $D$ be the interior of the ellipse.  For $z$ inside $D$ use its interior
harmonic measure $\omega_z$ on $\Gamma$; outside $D$ use exterior harmonic
measure of total mass one; on $\Gamma$ use $\delta_z$.
Let $W$ be the exterior inverse of
$\mathcal J(w)=\xi+cw+d/w$, where $c=(A+B)/2$ and $d=(A-B)/2$.  Define
$e(z)=0$ in $D$ and $e(z)=\log|W(z)|$ outside.  The factorization
$\mathcal J(w_1)-\mathcal J(w_2)=(w_1-w_2)(c-d/(w_1w_2))$ and the Green
identity imply
\[
 \int K_0(x,y)\,d\omega_z(y)\le K_0(x,z)+e(z),\qquad
 I(\omega_z,\omega_w)\le K_0(z,w)+e(z)+e(w),
\]
where $K_0=-\log|\cdot|$ and $I$ denotes its mutual energy.  The interior
deficit is the Green function of $D$; the exterior deficit is the
nonnegative disk Green function; the cross-boundary deficit is zero.
For $w=\rho e^{iu}$, $\rho\ge1$, the equipotential $|W(z)|=\rho$ is
the confocal ellipse
\[
 \Gamma_\rho:
 z=\xi+(c\rho+d/\rho)\cos u+i(c\rho-d/\rho)\sin u.
\]
It equals $\Gamma$ at $\rho=1$.  If $z=\mathcal J(w)$ and
$\zeta=\mathcal J(v)\in\Gamma$, $|v|=1$, the exact factorization gives
\[
 |z-\zeta|=|w-v|\left|c-\frac{d}{wv}\right|
 \ge(\rho-1)(c-d/\rho)\ge B(\rho-1).
\]
Taking the infimum over $\zeta$ and using $\log\rho\le\rho-1$ proves
\[
 e(z)\le\operatorname{dist}(z,\Gamma)/B\le\delta/B
 \qquad(z\notin D,\ z\in T_\delta).
\]

For $x\in[0,2]$, set
$\chi(z)=1-(\Re z-\xi)^2/A^2-(\Im z)^2/B^2$.  On $D$,
$|\nabla\chi|\le2/B$, while $\chi(x)>1/5$.  The level curve
$\chi=1/8$ is more than $1/32$ from $[0,2]$, so the Green function there is
at most $\log160<6$.  The barrier $48\chi$ on the collar gives
\[
 g_D(z,x)\le96\delta/B\qquad(z\in D\cap T_\delta,\ x\in[0,2]).
\]
Since $A_0$ has source mass three and $\lambda\Re z$ is harmonic, the
interior field-sweep loss for $A_0-\phi_q$ is at most $288\delta/B$.
The term $-V$ is favorable in this estimate: $V$ is subharmonic, with
distributional Laplacian $2\pi$ times the source measure on
$[-1/2,0]$, so its harmonic sweep is at least its interior value and the
$-V$ contribution can be dropped from an upper bound.
For exterior $z=\mathcal J(w)$, the swept mean coordinate is
$\xi+c/\overline w+d/w$, within $(A+B)\delta/B$ of $z$.
For the base-$2$ terms the source bound is
$[1/2+(A+B)\log2]\delta/B<3\delta/B$: $A_0$ changes by exactly
$-3e(z)$, and the remaining coordinate displacement is at most
$(A+B)\delta/B$.  The additional affine term is at most
$\lambda(A+B)\delta/B<6\delta/B$ because
$\lambda=\log(q/2)<2$ and $A+B<3$.

The logarithmic source in $\phi_q$ needs a separate collar estimate.  Put
\[
 V(z)=\int_0^{1/2}\log|z+s|\,ds.
\]
For $0<\delta\le1/64$ and an exterior point with
$\operatorname{dist}(z,\Gamma)\le\delta$, the explicit primitive
\[
 V(z)=\Re\bigl((z+1/2)\operatorname{Log}(z+1/2)
              -z\operatorname{Log}z-1/2\bigr)
\]
on either side of the source segment gives
\[
 \nabla V(z)=\operatorname{Log}(z+1/2)-\operatorname{Log}z.
\]
For $z=\xi+A\cos u+iB\sin u$ on $\Gamma$, minimizing
$|z|^2$ gives its minimum at $u=\pi$, namely
$\xi-A=751/5000$; the derivative has no other critical point because
$A\xi>|A^2-B^2|$.  The disk of this radius about $0$ lies inside the ellipse,
so every exterior point satisfies $|z|\ge751/5000$.  For the radial
comparison path, $h\le\delta/B$ and
\[
 |\mathcal J((1+h)e^{iu})-\mathcal J(e^{iu})|
 \le(c+d)h=Ah\le\frac AB\delta
 =\frac{12783}{11450}\delta<1.117\delta.
\]
This is about $1.116\delta$ and is less than $0.02$ when
$\delta\le1/64$.  Thus all points used in the estimate have
$\Re z\ge-751/5000-(12783/11450)/64>-1/4$.  Therefore
\[
 \frac{|z+1/2|}{|z|}
 \le1+\frac1{2|z|}
 \le1+\frac{2500}{751}
 =\frac{3251}{751}<4.33,
 \qquad
 |\arg(z+1/2)-\arg z|\le\pi.
\]
Since $|z+1/2|^2-|z|^2=\Re z+1/4>0$, the ratio is at least one, and
the modulus term in the difference of logarithms is at most $\log(4.33)$.
The angle difference is at most $\pi$ because the two points lie on the
same horizontal line.  Also $e^{3/2}>4.33$ by the first six terms of its
positive series, so $\log(4.33)<3/2$.  Thus, on both sides of the cut,
\[
 |\nabla V(z)|
 \le\sqrt{(3/2)^2+(22/7)^2}<3.72.
\]
At the contact point $z_0=\xi-A=-751/5000$ the classical gradient is not
defined, but these are the bounded one-sided traces and they give the
required Lipschitz bound across the smooth ellipse and the source cut.  The
Joukowski boundary parametrization has speed at most $4/3$, so its
pullback is $6$-Lipschitz.  If the exterior pole is $w=\rho e^{iu}$, write
$\rho=1+h$.  The elementary estimate
\[
 \int_{-\pi}^{\pi}|t|P_{1+h}(t)\,\frac{dt}{2\pi}
 \le \frac{3\pi h}{8}\log(1+4/h^2)
\]
is used for $0<h\le1$ and follows from
$1-\cos t\ge2t^2/\pi^2$.  In the present collar
$h\le\delta/B\le25/1832<1$, so this is the required range.  Since
$\operatorname{dist}(\mathcal J(w),\Gamma)\ge Bh$ and $\delta\le1/64$,
these bounds imply
\[
 \left|\int_\Gamma V(\zeta)\,d\omega_z(\zeta)-V(z)\right|
 \le20\delta\log(1/\delta).
\]
Thus the exterior field loss is at most
$9\delta/B+20\delta\log(1/\delta)$.  Combining this with the interior
bound gives a uniform field loss at most
$288\delta/B+20\delta\log(1/\delta)$.

The three interaction terms change by at most $3b$, where
$b=\int e\,d\nu\le|\nu|\delta/B$.  Since $|\nu|\le1$, their total loss is
 at most $3\delta/B$.  Consequently the full transfer cost is at most
$291\delta/B+20\delta\log(1/\delta)
 <300\delta+20\delta\log(1/\delta)$ because $B=229/200>1$.
Thus
\[
 \mathcal E_q(\kappa+\nu_\delta,\mu)
       \ge\alpha(\lambda)-300\delta-20\delta\log(1/\delta),
       \qquad \frac{291}{B}<300.
\]
The dual problem is posed on the printed support interval \([1/2,20]\).
The frequency-cutoff argument below first replaces the full frequency sum by
the range \(R\subset\{g+1,\ldots,20m\}\), with a positive multiplicative
tail factor.  Hence every finite configuration used in the variational bound
has its frequency measure in this compact interval; the omitted tail is
accounted for separately in \eqref{eq:p0-frequency-cutoff}.

For completeness, the unbounded frequency sum is reduced to a finite one
before counting configurations.  For $H\ge g+1$, define
\[
 E_{m,q}(H)=\sum_{S\subset I_m}\frac1{d!}
 \int_{(m\Gamma)^d}
 \sum_{\substack{R\subset\{g+1,\ldots,H\}\\|R|=m}}
 Q_{m,q}(S,R,Z)\prod_{z\in Z}\frac{|dz|}{2\pi}.
\]
The Cauchy--Binet expansion~\cite{Gantmacher} gives
$|G_{m,q}|\le E_{m,q}(\infty)$.
For fixed $S,Z$, its frequency-dependent
factor is
\[
 f_q(R)=
 \frac{\prod_{r\in R}w_{r,q}\,\Delta(R)^2
       \prod_{r\in R,z\in Z}|r+z|}
      {\prod_{r\in R,i\in I_m}(r+i)
       \prod_{r\in R,j\in J_m\setminus S}(r+j)}.
\]
If $r=\max R>20m$, replace $r$ by the least missing
$s\in\{g+1,\ldots,g+m\}$, so $s\le3m/2$.  The exact ratio of
the frequency weights, both frequency Vandermonde factors, and all
contour-frequency factors gives
\[
\frac{f_q(R)}{f_q(R')}\le
      12500^m\frac{m^2}{4r^2}w_{r,2}.
\]
For the constant, at most two integers lie at each distance from $s$, so
\[
 \prod_{u\in R\setminus\{r\}}|s-u|
 \ge d_m:=
 \left\lfloor\frac{m-1}{2}\right\rfloor!
 \left\lceil\frac{m-1}{2}\right\rceil!
 \ge e^{-1}\left(\frac{m}{2e}\right)^{m-1}.
\]
Cancel the common factors in the two Cauchy determinants.  The two
frequency Vandermonde ratios contribute the squared products of distances
from $r$ and $s$; each contour-frequency factor is bounded by the
corresponding ratio involving $r+g$ and $s+g$, and the fixed-grid factors
are bounded by $5^m$.  The explicit denominator bounds
$s+i\le11m/4$, $s+j\le7m/2$, and ellipse bounds
$\Re z>-m/6$, $|z|<5m/2$ give
\[
 \frac{f_q(R)}{f_q(R')}
 \le \frac{w_{r,q}}{w_{s,q}}
       \left(\frac{77}{8}\right)^m
       \left(\frac{189}{16}\right)^{m-|S|}
       \frac{m^{2m}}{r^2d_m^2}.
\]
For $r>s$, $w_{r,q}/w_{s,q}\le w_{r,2}/w_{s,2}$ and
$w_{s,2}\ge2^{-3m/2}$ because $s\le3m/2$ and
$\binom{s-1}{g}\ge1$.  Substituting this and $e<3$ into the last
display gives the asserted $12500^m m^2w_{r,2}/(4r^2)$ bound.  This is the
full constant calculation; it does not require an estimate for $w_{s,q}$.
The exact base comparison is
\[
 \frac{w_{r,q}/w_{s,q}}{w_{r,2}/w_{s,2}}
       =\left(\frac2q\right)^{r-s}\le1.
\]
\begin{proof}[Frequency cutoff]\label{eq:p0-frequency-cutoff}
\[
 \vartheta_m=\frac{m}{1600}12500^m\sum_{r>20m}w_{r,2}.
\]
Indeed, the replacement map has at most $m$ preimages for each pair
$(R',r)$, and $m^2/(4r^2)<1/1600$ when $r>20m$.  For $t=15/8$,
\[
 \sum_{r>20m}w_{r,2}
 \le(8/15)^{20m}\sum_{r\ge g+1}w_{r,2}(15/8)^r
 =(8/15)^{20m}15^{g+1}.
\]
Since $g=m/2$ and
$12500\sqrt{15}<50000$, this gives
\[
 \vartheta_m\le\frac{15m}{1600}
       \left[50000(8/15)^{20}\right]^m
 <\frac3{320}2^{-m}<\frac12.
\]
The strict inequality uses $200000\cdot8^{20}<15^{20}$ and
$m2^{-m}\le1$.  If $Z_M$ denotes the finite frequency sum through $M$,
the replacement map gives $Z_M-Z_{20m}\le \vartheta_mZ_M$.  Hence
$Z_M\le(1-\vartheta_m)^{-1}Z_{20m}\le(1+2\vartheta_m)Z_{20m}$.
\[
 2\vartheta_m<\frac3{160}2^{-m}.
\]
Monotone convergence as $M\to\infty$ now proves
\[
 E_m(\infty)\le\left(1+\frac3{160}2^{-m}\right)E_m(20m).
\]
\end{proof}
There are $2^m$ pinned subsets and at most $\binom{39m/2}{m}$ retained
frequency sets.  The normalized contour length is at most $mA$, hence
$(mA)^d/d!\le e^{mA}$.  Using $c_*<49.979$ and
$c_*+\log39+1+A<56$, the finite configuration bound, the tube transfer
with $\delta=m^{-2}$, and the tail factor imply the following.

\begin{lemma}[Uniform smoothing bound]\label{lem:p0-smoothing-transfer}
For every integer $q\in\{2,\ldots,11\}$ and every $m\in4\mathbb N$,
 $m\ge8$,
\begin{equation}\label{eq:p0-q-uniform-det}
 \log|G_{m,q}(\Li_2(1/q))|
 \le-\alpha(\log(q/2))m^2+10m\log m+58m+301+40\log m.
\end{equation}
The frequency quantifier before replacement is the full set
$R\subset\{g+1,g+2,\ldots\}$ of cardinality $m$.  The constant $58$ is
uniform in $q$; the constant $301$ includes the tube cost $300$, the
real-cell remainder $1/3$, and the exponentially small frequency-tail
factor.  The additional $40\log m$ is the explicit logarithmic collar cost.
\end{lemma}

\begin{proof}
For each finite configuration the cell measures above are admissible in the
$m^{-2}$ tube problem.  Tube transfer bounds their energy below by
$\alpha(\lambda)-300m^{-2}-40m^{-2}\log m$.  Inserting this into the pointwise determinant
envelope gives a $300+40\log m$ additive loss.  Summing the $2^m$ pinned choices,
the $\binom{39m/2}{m}$ retained frequency choices, and the contour integral
 adds at most $(\log39+1+A)m$ beyond $c_*m$.  The first contribution is
strictly less than $56m$ by the displayed $c_*+\log39+1+A<56$; the extra
diagonal factor is strictly less than $e^{2m}$.  Hence the coefficient of
$m$ is strictly less than $56+2=58$.  The cell remainder is $1/3$.  For
$m\ge8$ the tail factor obeys
\[
 \log\left(1+\frac3{160}2^{-m}\right)
 \le\frac3{160}2^{-m}\le\frac3{40960},
\]
so $300+1/3+3/40960<301$.  Taking the
monotone limit of the finite frequency cutoffs proves the estimate for the
full half-line.
\end{proof}

\subsection{Endpoint certificates and the interval comparison}

The exact reference masses and thresholds are printed in
Tables~\ref{tab:q2data} and~\ref{tab:q11data}.  Their directed dual checks
prove
\begin{equation}\label{eq:p0-q-endpoint-bounds}
 \alpha(0)\ge\frac{248035}{125000},\qquad
 \alpha(\log(11/2))\ge\frac{77597}{100000}.
\end{equation}
These are lower bounds for the full continuum infimum, not empirical
minimizers.  In each certificate the contour inequality is checked on the
whole ellipse; the pinned positive-part integral is checked on its full
support; the frequency positive-part integral is checked through $81/10$.
The endpoint inequalities at $81/10$ are part of the directed integrals
recomputed from the printed measures: they give $W_\mu(81/10)>31/10$
for $q=2$ and $W_\mu(81/10)>56/5$ for $q=11$, while the respective thresholds
are $1427/500$ and $1217/200$.

The infinite tail has a direct bound that uses only the masses and supports
of these two references.  Both $\eta_0$ have mass one and support in
$E\cup\Gamma$, whose real part is at least $-751/5000$; both $\mu_0$ have
mass one and support in $[1/2,9/2]$.  Also $B_0'$ is the Cauchy transform
of a positive measure of mass three on $[0,2]$.  Differentiating the
definition of $W_\mu$ therefore gives, for $t\ge81/10$,
\begin{align*}
 W_\mu'(t)
 &=B_0'(t)+\log q-\log\frac{t}{t-1/2}
   -2\int\frac{d\mu_0(s)}{t-s}
   -\int\Re\frac{d\eta_0(z)}{t+z}\\
 &\ge H(t):=\frac3{t+2}+\log q-\log\frac{t}{t-1/2}
       -\frac2{t-9/2}-\frac1{t-751/5000}.
\end{align*}
For $t\ge81/10$, differentiation yields
\[
H'(t)=-\frac3{(t+2)^2}+\frac1{2t(t-1/2)}
       +\frac2{(t-9/2)^2}+\frac1{(t-751/5000)^2}>0,
\]
because the last two positive terms exceed $3/t^2$, while
$3/(t+2)^2<3/t^2$.  Thus $H(t)\ge H(81/10)$.  The elementary bounds
$\log q\ge\log2>2/3$, $\log(81/76)<1/13$, and
$5000/39749<13/100$ give
\[
W_\mu'(t)\ge H(81/10)>
\frac{30}{101}+\frac{69}{100}-\frac1{13}-\frac59-\frac{13}{100}
=\frac{66338}{295425}>\frac{31}{684}.
\]
Here we used the elementary rational-series bounds
$\log2>69/100$, $\log11>23/10$, $\log(81/76)<1/13$, and
$5000/39749<13/100$; the $q=11$ lower bound is larger than $1$.
Consequently $W_\mu'(t)>31/684>1/50$ for $q=2$ and is $>1$ for
$q=11$.  Hence $W_\mu(t)>h_\mu$ on the entire omitted ray, so $T$ has no
omitted tail.  The directed endpoint records are summarized in
Table~\ref{tab:p0-q-bounds}; the exact mass lists are the two preceding
reference tables.

For auditability, the scalar inputs and the outward rational penalties used
by the two records are listed here (the displayed finite decimals are
outward lower bounds for $Q_{\log}$):
\begin{table}[ht]
\centering
\caption{Scalar data in the two directed endpoint certificates.}
\label{tab:p0-q-bounds}
\renewcommand{\arraystretch}{1.12}
 {\normalsize
\begin{tabular}{@{}c l@{}}
\toprule
$q$ & exact thresholds and directed penalty bounds\\
\midrule
$2$ & $h_\eta=-3347002/1000000$, $h_\mu=1427/500$;\\
    & $Q_{\log}>-0.425508$,\quad
      $P\le0.087950$,\\
    & $T\le0.122962$;\\[2pt]
$11$ & $h_\eta=-3359/500$, $h_\mu=1217/200$;\\
    & $Q_{\log}>-0.635537$,\quad $P\le0.000001$,\\
    & $T\le1.069189$.\\
\bottomrule
\end{tabular}}
\end{table}

The arithmetic rate used for the positive bases is
\[
 \tau_+(\lambda)=K+\frac1{16}\log(q-1)-\frac{13}{16}\log q,
                    \qquad q=2e^\lambda,
\]
where the arithmetic estimate gives $K\le1271001/500000$.  Differentiation
with respect to $\lambda$ gives
\[
 \tau_+''(\lambda)=-\frac{q}{16(q-1)^2}<0.
\]
Its tangent at $q_0=63/25$ has exact slope
\[
 \tau_+'(\log(q_0/2))=-\frac{431}{608}.
\]
Concavity places $\tau_+$ below this tangent on the whole interval.  Directed
logarithm evaluation with the stated rational upper bound for $K$ gives the
endpoint comparison:
\begin{center}
\renewcommand{\arraystretch}{1.12}
 {\normalsize
\begin{tabular}{@{}c c c c@{}}
\toprule
$q$ & directed $\alpha$ lower bound & tangent upper bound for $\tau_+$
     & directed gap lower bound\\
\midrule
$2$ & $248035/125000$ & $1.981041880111$ & $>3/1000$\\
$11$ & $77597/100000$ & $0.772577360778$ & $>3/1000$\\
\bottomrule
\end{tabular}}
\end{center}
The displayed decimal upper bounds are outward roundings of the directed
intervals; the corresponding directed gaps are respectively
$0.00323811988930149\ldots$ and $0.00339263922279043\ldots$.

Let $L(\lambda)$ be the chord joining the two lower bounds in
\eqref{eq:p0-q-endpoint-bounds}.  By concavity of $\alpha$,
$\alpha(\lambda)\ge L(\lambda)$.  Let $T(\lambda)$ be the tangent to
$\tau_+$ at $q_0=63/25$.  The two directed endpoint checks show
$L-T>3/1000$ at both ends.  Since $L-T$ is affine in $\lambda$, it is
greater than $3/1000$ throughout
$0\le\lambda\le\log(11/2)$.  Also $\tau_+(\lambda)\le T(\lambda)$.
We have proved
\begin{equation}\label{eq:p0-q-positive-margin}
 \alpha(\lambda)-\tau_+(\lambda)>\frac3{1000}
                 \quad(0\le\lambda\le\log(11/2)).
\end{equation}
The determinant estimate \eqref{eq:p0-q-uniform-det} is a separate
smoothing estimate; the fixed-base integerizer and nonvanishing sequence are
applied separately at each integer $q$.  Thus this section uses
no interpolation of integerizers, seed primes, or endpoint reference
measures.
